\section{Quantidade de movimento angular do sistema}
\label{sec_dinAcoplada_quantidade_movimento_angular}

Para fins de análise física e verificação da conservação da quantidade de movimento angular,
introduz-se o cálculo explícito da quantidade de movimento angular total do sistema
espaçonave--\gls{mr} em relação à origem do referencial $\{B\}$.
Na formulação de velocidades espaciais adotada na Seção~\ref{sec_dinAcoplada_energia_cinetica},
a velocidade espacial do corpo $i$ (espaçonave ou elo do \gls{mr}) em $\{B\}$ é escrita como:

\begin{equation}
\vec v_{i,\text{sp}}
=
\begin{bmatrix}
{}^{B}\!\vec v_{G_i} \\[2pt]
{}^{B}\!\vec\omega_i
\end{bmatrix}
=
J_{\text{sp},i}(\vec y\,)\,\dot{\vec y\,}.
\label{eq_velocidade_espacial_corpo_i}
\end{equation}


Em que ${}^{B}\!\vec v_{G_i}$ e ${}^{B}\!\vec\omega_i$ são, respectivamente,
a velocidade linear do \gls{cdm} e a velocidade angular do corpo $i$
expressas em $\{B\}$, e $J_{\text{sp},i}(\vec y\,)$ é o jacobiano espacial
avaliado no \gls{cdm} do corpo $i$, construído a partir dos graus de
liberdade de atitude da espaçonave e das juntas do manipulador
\cite{featherstone2008}.

Para cada corpo $i$, o operador de inércia espacial
$M_{i,\text{spt}}$ associa a velocidade espacial à quantidade de movimento espacial
\cite{featherstone2008}:

\begin{equation}
\vec h_{i,\text{sp}}
=
M_{i,\text{spt}}\,\vec v_{i,\text{sp}}.
\label{eq_quantidade_movimento_espacial_corpo_i}
\end{equation}


Em que $\vec h_{i,\text{sp}}$ agrupa em um único vetor a quantidade de movimento linear e a quantidade de movimento angular do corpo $i$ em relação à origem de $\{B\}$.
A parte angular de $\vec h_{i,\text{sp}}$ é obtida selecionando-se as
três últimas componentes, por meio da matriz:

\begin{equation}
S_\omega
=
\begin{bmatrix}
0 & 0 & 0 & 1 & 0 & 0 \\
0 & 0 & 0 & 0 & 1 & 0 \\
0 & 0 & 0 & 0 & 0 & 1
\end{bmatrix},
\qquad
\vec h_i
=
S_\omega\,\vec h_{i,\text{sp}}.
\label{eq_selecao_parte_angular}
\end{equation}

Em que $\vec h_i$ é o vetor de quantidade de movimento angular do corpo $i$ expresso em $\{B\}$,
em analogia à decomposição da quantidade de movimento espacial em componentes linear e angular discutida em
\cite{featherstone2008,lynch2017modern}.

Neste trabalho, a espaçonave é tratada como um corpo rígido adicional com operador de inércia espacial
$M_{\text{espaç},\text{spt}}$ expresso em $\{B\}$,
cujos parâmetros de massa, \gls{cdm} e matriz de inércia já incluem
a contribuição estrutural da espaçonave.
Os elos do manipulador são indexados por $i=1,\dots,n$, de modo que o
quantidade de movimento angular total do sistema espaçonave e manipulador em $\{B\}$ é escrito como:

\begin{equation}
\vec h_{\text{tot}}
=
\sum_{i=0}^{n} \vec h_i
=
\sum_{i=0}^{n}
S_\omega\,M_{i,\text{spt}}\,
J_{\text{sp},i}(\vec y\,)\,\dot{\vec y\,}.
\label{eq_quantidade_mov_angular_total}
\end{equation}


Essa expressão define a quantidade de movimento espacial total como a soma dos quantidade de movimentos de cada corpo rígido
em sistemas de espaçonave flutuante \cite{featherstone2008} e com formulações de sistemas espaçonave--manipulador
que utilizam a conservação de quantidade de movimento para descrever a dinâmica em
microgravidade.
Na ausência de torques externos, a quantidade de movimento angular total do sistema em relação
ao ponto de referência adotado (origem de $\{B\}$) é conservada, de modo que $\vec h_{\text{tot}}$
deve permanecer constante no tempo \cite{wie2008space,featherstone2008}.
